% ============================================================
%  1.1 – Xác định dự án  (0.5 điểm)
%  Trạng thái: ĐÃ HOÀN THÀNH
% ============================================================

\subsection{Bối cảnh và vấn đề}

Trong môi trường đại học, các hoạt động hỗ trợ học tập, nghiên cứu và tổ chức phong trào
diễn ra thường xuyên và liên tục. Sinh viên có nhu cầu tìm gia sư các môn chuyên ngành, tìm
trợ giảng, tìm người hỗ trợ phòng thí nghiệm, cộng tác viên tổ chức sự kiện, hoặc tham gia
các hoạt động công tác xã hội do Đoàn--Hội phát động. Đồng thời, giảng viên, phòng ban và
các câu lạc bộ trong trường cũng thường xuyên cần tuyển sinh viên tham gia công việc ngắn
hạn hoặc bán thời gian.

Mặc dù các hoạt động này diễn ra với tần suất cao, quá trình kết nối giữa bên có nhu cầu và
bên cung ứng chủ yếu thực hiện qua các kênh \textbf{phi chính thức} như Facebook Group nội bộ,
Zalo nhóm lớp, email thông báo hoặc truyền miệng. Qua quan sát thực tế, nhóm xác định
\textbf{ba vấn đề cốt lõi}:

\begin{enumerate}
  \item \textbf{Thông tin phân tán, thiếu hệ thống.}
    Bài đăng trên Facebook nhanh chóng bị trôi khỏi dòng thời gian; không có cơ sở dữ liệu
    tập trung cho phép lọc theo loại công việc, đơn vị tuyển dụng hoặc kỹ năng yêu cầu.

  \item \textbf{Thiếu cơ chế đánh giá và minh bạch năng lực.}
    Không có hồ sơ kỹ năng chuẩn hóa, không có lịch sử đánh giá sau hoàn thành công việc;
    tuyển chọn chủ yếu dựa trên sự quen biết.

  \item \textbf{Quy trình giao việc không được số hóa.}
    Không có bước xác nhận chính thức trạng thái ``đang thực hiện'' hay ``đã hoàn thành'';
    thanh toán hoàn toàn dựa trên niềm tin cá nhân, không có cơ chế trung gian.
\end{enumerate}

\subsection{Nghiên cứu thị trường -- Phân tích hệ sinh thái nội bộ}

Nhóm khảo sát bốn kênh hiện đang được sinh viên, giảng viên và tổ chức trong trường sử dụng:

\begin{table}[H]
\centering
\caption{So sánh các kênh kết nối hiện tại trong môi trường đại học}
\begin{tabularx}{\textwidth}{|p{2.4cm}|X|X|}
\hline
 \textbf{Kênh} & \textbf{Ưu điểm} & \textbf{Hạn chế} \\
\hline
Facebook Group & Độ phủ cao, quen thuộc & Thông tin trôi nhanh; không có cấu trúc; không theo dõi trạng thái \\
\hline
Zalo Group & Tức thời cao & Thông tin bị mất trong nhóm đông; không phân loại hay lưu trữ \\
\hline
Email trường & Tính chính thống cao & Giao tiếp một chiều; không có danh sách vị trí mở; không có cơ chế tracking \\
\hline
Google Form & Chuẩn hóa đầu vào & Chỉ thu thập dữ liệu; không xây dựng hồ sơ năng lực lâu dài \\
\hline
\end{tabularx}
\end{table}

\textbf{Kết luận:} Không có kênh nào trong số trên tạo ra quy trình giao dịch khép kín với
đánh giá, xác nhận hoàn thành và lịch sử năng lực tích lũy.

\subsection{Ý tưởng dự án và mô hình}

Xuất phát từ khoảng trống trên, nhóm đề xuất dự án \textbf{``UniJob'' (Uni Task Hub)} -- một
nền tảng thương mại điện tử nội bộ dành riêng cho cộng đồng trong trường đại học.

\begin{itemize}
  \item \textbf{Mô hình:} C2C + B2C nội bộ (sinh viên, giảng viên, Đoàn--Hội, CLB, phòng ban).
  \item \textbf{Quy trình:} Đăng task $\rightarrow$ Ứng tuyển $\rightarrow$ Xác nhận nhận việc
    $\rightarrow$ Thực hiện $\rightarrow$ Xác nhận hoàn thành $\rightarrow$ Đánh giá.
  \item \textbf{Điểm khác biệt:} Xác thực qua email trường học (\texttt{@hcmut.edu.vn}), giảm
    rủi ro ẩn danh, tăng độ tin cậy nội bộ.
\end{itemize}

\subsection{Xác định khách hàng}

Nền tảng phục vụ hai nhóm người dùng song song trong cùng một hệ sinh thái:

\begin{itemize}
  \item \textbf{Bên đăng task:} Giảng viên cần trợ giảng / hỗ trợ nghiên cứu; phòng ban cần
    nhập liệu / trực văn phòng; CLB cần cộng tác viên sự kiện; sinh viên cần gia sư.
    \textit{Nhu cầu chung:} tuyển người nhanh, tin cậy, theo dõi tiến độ.

  \item \textbf{Bên nhận task:} Sinh viên muốn tích lũy kinh nghiệm, xây dựng hồ sơ hoạt động,
    có thêm thu nhập hoặc điểm rèn luyện.
    \textit{Nhu cầu chung:} môi trường minh bạch, có thể chứng minh năng lực qua lịch sử
    và đánh giá.
\end{itemize}

\subsection{Tính đổi mới và tính khả thi}

\textbf{Tính đổi mới:} Không nằm ở công nghệ phức tạp mà ở việc tái cấu trúc hệ sinh thái
task sẵn có thành một ``campus marketplace'' chuyên biệt với quy trình số hóa đầy đủ.

\textbf{Tính khả thi kỹ thuật:} Website xây dựng với React 19 + Firebase (Auth, Firestore,
Storage, Hosting); không yêu cầu logistics hay chuỗi cung ứng phức tạp; chi phí vận hành thấp.

\textbf{Tính khả thi vận hành:} Phạm vi ban đầu giới hạn trong một trường; xác thực qua email
trường giúp hạn chế tài khoản giả mạo; quản lý nội dung đơn giản.

\textbf{Tính khả thi pháp lý:} Nền tảng hoạt động như hệ thống kết nối dịch vụ nội bộ, không
kinh doanh hàng hóa cấm, không vi phạm quy định về TMĐT Việt Nam.

\subsection{Kết luận}

Dự án UniJob được đề xuất nhằm số hóa và chuẩn hóa quy trình giao--nhận task nội bộ, xây
dựng marketplace chuyên biệt trong phạm vi trường đại học. Với phạm vi triển khai hẹp, chi
phí vận hành thấp và nhu cầu thực tế cao, dự án có tính khả thi đáng kể cả về kỹ thuật, vận
hành và chiến lược phát triển dài hạn.
